\section{Introduction}
\label{sec:intro}

Hand gestures are a natural way to command a mobile robot in places where
speech is unreliable, such as a busy lobby, a laboratory or a caf\'e.
Legged robots are increasingly deployed in exactly these spaces, and
off-the-shelf pipelines now detect and classify hand gestures in real time on
modest hardware~\cite{zhang2020mediapipehands,kapitanov2024hagrid}. Most
gesture interfaces, however, assume a single user standing in front of the
robot: a hand is detected, classified and mapped to a velocity command. In a
crowd this assumption breaks. Several people are in view, many of them
gesture to each other, and the robot cannot tell a command from a
conversation. A bystander's open palm stops the robot; a fist meant for a
friend sends it walking.

The question a gesture interface must answer in a shared space is therefore
not only \emph{which} gesture is shown, but \emph{who} is showing it and
\emph{whether it is addressed to the robot}. We formulate three requirements:
\begin{enumerate}
\item[R1] gestures from different people are recognised and attributed to
  the right person, even when they occur at the same time;
\item[R2] only gestures addressed to the robot are executed;
\item[R3] the same person is ignored when they turn and gesture to someone
  else.
\end{enumerate}

Humans resolve this with nonverbal cues: people who address someone face
them, look at them and stand within a social distance~\cite{hall1966hidden,
kendon1990conducting}. These cues are easy to state as linguistic rules
(``\emph{if} a person is facing the robot \emph{and} is at a social distance
\emph{then} they are engaged''), which makes fuzzy logic a natural model.
The difficulty is that the cues are measured from a moving robot: depth
readings have holes and edge errors on people, gaze is only resolvable at
short range, and orientation estimates degrade with occlusion and motion
blur. A
Type-1 fuzzy system treats every reading as exact. An Interval Type-2 (IT2)
system can represent how much each reading should be trusted, but in most
applications its footprint of uncertainty (FOU) is fixed at design time.

We propose an engagement gate in which the FOU of every input is driven at
run time by the measured uncertainty of that input for that person. When a
person is clearly visible the IT2 system behaves like its Type-1 counterpart;
when they are far away, partly occluded or turning, the engagement interval
widens and the decision, taken on its lower bound, becomes conservative. The
contributions of this paper are:
\begin{itemize}
\item a per-person perception pipeline for a quadruped robot that attributes
  gestures to individuals and estimates distance, approach speed,
  orientation and gaze, each with an explicit uncertainty, from a single
  RGB-D camera mounted on the robot;
\item a two-layer fuzzy engagement gate for gesture commands, and an IT2
  variant whose FOU is set per person and per frame from measured perception
  uncertainty, reducing exactly to the Type-1 system when uncertainty is
  zero;
\item an evaluation on a real Unitree Go2 in stationary and walking
  scenarios with several people, comparing gesture-only, nearest-person,
  gaze-only, learned, Type-1 and IT2 policies on identical inputs, with an
  analysis of where measured uncertainty helps.
\end{itemize}

The remainder of the paper is organised as follows. Section~\ref{sec:background}
reviews IT2 fuzzy systems, Section~\ref{sec:related} discusses related work,
Section~\ref{sec:method} describes the method, Section~\ref{sec:experiments}
the experimental setup, Section~\ref{sec:results} the results, and
Section~\ref{sec:discussion} limitations and future work.
